\section{The rank of a homogeneous depth-five circuit}
\label{sec:circuits}

We now bound the rank measure in terms of circuit size.  The argument
is uniform over the choice of the $k$ matrix layers assigned to $V$.
Its first step bounds the rank of a product using only the degrees of
its factors.  We then expand one middle-layer sum when its degree is
large, exposing its lower products of linear forms.

\begin{proposition}
\label{prop:circuit}
There are absolute constants $C>0$ and $n_0$ with the following property.
For $n\ge n_0$, use the parameters in~\eqref{eq:parameters} and any
partition of the matrix layers with $k$ layers in $V$ and $m$ layers in
$U$.  Let $K$ be a field, and let $g\in K[V,U]$ be a homogeneous
degree-$n$ polynomial computed by a syntactically homogeneous
$\Sigma\Pi\Sigma\Pi\Sigma$ circuit of size $S$.  Then
\[
 R(g)\le 2S^2D\exp(C\sqrt n)\,n^{-\sqrt n/100}.
\]
\end{proposition}

All thresholds in this section are absorbed into a single absolute
$n_0$, independent of the field, the circuit, and the partition.
Only the parameter estimates of Lemma~\ref{lem:parameters} will be
used.

\subsection{Products and intermediate dimensions}

For a positive integer $e$, define
\[
 H_e=\sum_{i=0}^{e}q^{|i-\lambda e|}.
\]
These quantities account for all choices of bidegree in a product.

The quantities $H_e$ use the same integer-degree discrepancies as the
residue method of Amireddy, Garg, Kayal, Saha, and Thankey
\cite[Definition~2.1 and Lemma~3.1]{AmireddyGargKayalSahaThankey2022Full}.
In the present argument they control the dimensions of the intermediate
homogeneous spaces through which the differentiation--multiplication
operator factors.

\begin{lemma}
\label{lem:product-rank}
Let $K$ be a field, and let $Q_1,\ldots,Q_r\in K[V,U]$ be homogeneous of respective positive
degrees $e_1,\ldots,e_r$, with $\sum_{j=1}^r e_j=n$.  Then
\[
 R\left(\prod_{j=1}^r Q_j\right)
 \le 2D\prod_{j=1}^r H_{e_j}.
\]
\end{lemma}

\begin{proof}
Decompose each factor into its bidegree components.  The component of
their product that defines the rank measure is
\[
 \left(\prod_{j=1}^r Q_j\right)_{[k,m]}
 =\sum_{\substack{0\le i_j\le e_j\ (1\le j\le r)\\
                   \sum_j i_j=k}}
   \prod_{j=1}^r (Q_j)_{[i_j,e_j-i_j]}.
\]
Fix an assignment in this sum, and put
$\delta_j=i_j-\lambda e_j$.  Since $\lambda n=k$, we have
$\sum_j\delta_j=0$.  The operators obtained by substituting
$\partial_V$ and $U$ commute, so we may apply first the factors with
$\delta_j>0$.  Write their total bidegree as
\[
 I=\sum_{\delta_j>0}i_j,
 \qquad J=\sum_{\delta_j>0}(e_j-i_j).
\]
The resulting intermediate space is
$\HH_V(a-I)\otimes\HH_U(b+J)$.  It is well defined because
$I\le k\le a$; also $J\le m$.  Its dimension bounds the rank of this
contribution, including when the set of positive discrepancies is
empty.

The successive ratios of binomial coefficients give
\begin{align*}
 \frac{h(v,a-I)}{h(v,a)}
 &=\prod_{t=0}^{I-1}\frac{a-t}{a+v-1-t}
 \le\left(\frac{a}{a+v-1}\right)^I,\\
 \frac{h(u,b+J)}{h(u,b)}
 &=\prod_{t=1}^{J}\frac{b+u+t-1}{b+t}
 \le\left(\frac{b+u}{b}\right)^J
 =\beta^{-J}.
\end{align*}
The first inequality follows because $x/(x+v-1)$ is nondecreasing for
$x>0$.  Since $\beta\ge\alpha^\rho$, these inequalities imply
\begin{equation}
\label{eq:intermediate-dimension}
 \frac{h(v,a-I)h(u,b+J)}{D}
 \le \alpha^{I-\rho J}
       \left(1-\frac1{a+v}\right)^{-I}
 \le 2\alpha^{I-\rho J}.
\end{equation}
The last bound holds uniformly for $I\le n$, because
$a+v=\Theta(n^3)$.  Empty products in the ratio formulas equal 1.

Using $\lambda=\rho/(1+\rho)$ and $\sum_j\delta_j=0$, we obtain
\[
 I-\rho J
 =(1+\rho)\sum_{\delta_j>0}\delta_j
 =\frac{1+\rho}{2}\sum_j|\delta_j|.
\]
As $q=\alpha^{(1+\rho)/2}$, the rank of the fixed-assignment
contribution is therefore at most
\[
 2D\prod_j q^{|i_j-\lambda e_j|}.
\]
Rank subadditivity permits summing these bounds.  Dropping the
constraint $\sum_j i_j=k$ increases the resulting nonnegative sum
and makes it factor as $2D\prod_jH_{e_j}$.  Zero bidegree components
contribute rank zero throughout.
\end{proof}

\subsection{The cost of integer degrees}

The product estimate is effective because $\lambda e$ is usually
separated from the integers.  Small even degrees require particular
care: their separation can tend to zero, so a constant loss for each
factor would be too large.  We bound their combined error instead.

Put
\[
 d_e=\dist(\lambda e,\Z),\qquad t_0=\frac{n}{4s}.
\]
\begin{lemma}
\label{lem:degree-estimates}
For every integer $e\ge1$,
\begin{equation}
\label{eq:nearest-integer-bound}
 H_e\le q^{d_e}\bigl(1+4q^{1-2d_e}\bigr)\le5.
\end{equation}
If $1\le e<t_0$, then
\begin{equation}
\label{eq:low-degree-distances}
 d_e=
 \begin{cases}
   se/(2n),&e\text{ even},\\
   1/2-se/(2n),&e\text{ odd},
 \end{cases}
 \qquad d_e\ge\frac{se}{2n}.
\end{equation}
In this range, odd degrees satisfy $H_e\le q^{d_e/2}$, whereas even
degrees satisfy
\begin{equation}
\label{eq:even-degree-bound}
 H_e\le q^{d_e}\exp\bigl(4q^{1-se/n}\bigr).
\end{equation}
There is an absolute constant $C_1$ such that, for every finite list of
positive integers $e_1,\ldots,e_r$ with $\sum_j e_j\le n$,
\begin{equation}
\label{eq:even-corrections}
 \sum_{\substack{j:\ e_j<t_0\\ e_j\text{ even}}}
 q^{1-se_j/n}\le C_1\sqrt n.
\end{equation}
\end{lemma}

\begin{proof}
Enlarge the sum defining $H_e$ to all integers.  Apart from one nearest
integer to $\lambda e$, the terms have total at most
$2q^{1-d_e}/(1-q)$.  This remains valid when the nearest integer is
not unique.  Since $q\le1/2$, it gives the first inequality
in~\eqref{eq:nearest-integer-bound}; the second follows from
$0\le d_e\le1/2$ and $0<q<1$.

For $e<t_0$, write
\[
 \lambda e=\frac e2-\frac{se}{2n},
 \qquad 0<\frac{se}{2n}<\frac18.
\]
The nearest integer is consequently $e/2$ for even $e$, and
$(e-1)/2$ for odd $e$.  This proves
\eqref{eq:low-degree-distances}.  For odd $e$ we have $d_e\ge3/8$.
By $q\le n^{-2/5}$,
\[
 q^{-d_e/2}\ge n^{d_e/5}\ge n^{3/40}\ge5
\]
after increasing $n_0$.  Thus
$H_e\le5q^{d_e}\le q^{d_e/2}$.  For even $e$, substitute
$2d_e=se/n$ in~\eqref{eq:nearest-integer-bound} and use
$1+x\le\exp(x)$ to obtain~\eqref{eq:even-degree-bound}.

It remains to sum these even-degree errors.  Increase $n_0$ so that
$t_0>2$, and for real $z\in[2,t_0]$ define
\[
 f(z)=\frac{q^{1-sz/n}}{z}.
\]
Its logarithm is convex, since $(\log f)''(z)=1/z^2>0$.
Hence $f(z)\le\max\{f(2),f(t_0)\}$.  The sharp estimate for $q$ in
Lemma~\ref{lem:parameters} gives
\[
 f(2)
 =\frac12 n^{-1/2}
   \exp\bigl(O(s\log n/n+1/a)\bigr)
 =O(n^{-1/2}).
\]
At the other endpoint, $s\le\sqrt n$ and $q\le2n^{-1/2}$ give
\[
 f(t_0)=\frac{4s}{n}q^{3/4}=O(n^{-7/8}).
\]
Therefore $q^{1-se/n}\le C_1e/\sqrt n$ uniformly over the low even
degrees.  Summing and using $\sum_j e_j\le n$ proves
\eqref{eq:even-corrections}.
\end{proof}

These estimates control a product of factors below $t_0$ regardless
of their coefficients.  A factor of degree at least $t_0$ will instead
supply many linear factors through its bottom-layer expansion.

\subsection{From gates to products}

\begin{proof}[Proof of Proposition~\ref{prop:circuit}]
If $g=0$, the conclusion is immediate.  Otherwise its formal output
degree is $n$: induction on the gates shows that every nonzero
polynomial at a formally homogeneous degree-$d$ gate is homogeneous
of degree $d$.

We first record exactly what the circuit layers provide.  A bottom
sum gate computes either a homogeneous linear form or a scalar,
because all its leaf inputs have the same formal degree.  A lower
product of formal degree $e>0$ therefore computes a scalar times a
product of exactly $e$ homogeneous linear-form occurrences.
Degree-zero factors are absorbed into that scalar.  A zero factor
makes the whole product zero and its contribution can be discarded.
A gate computing zero need not be assigned a different formal degree.

At either a middle sum or the output sum, repeated wires from the
same predecessor gate can be combined by adding their scalar
coefficients.  There are at most $S$ distinct predecessors.  Thus
$g$ is a sum of at most $S$ scalar multiples of products
\begin{equation}
\label{eq:circuit-product-form}
 G=\prod_{j=1}^{r}Q_j,
 \qquad \deg Q_j=e_j\ge1,\qquad \sum_j e_j=n,
\end{equation}
and each $Q_j$ is a sum of at most $S$ scalar multiples of products
of $e_j$ homogeneous linear forms.  Zero upper terms have been
discarded, and degree-zero factors in the remaining terms have been
absorbed into their coefficients.  If empty sums or products are allowed, they
contribute only zero or a scalar and are handled in the same way.

In~\eqref{eq:circuit-product-form}, inputs to a product are counted
with multiplicity: repeated use of a gate gives repeated factors.
Their positive degrees add to $n$, so $r\le n$.  Shared subcircuits
do not affect the bounds on the number of distinct additive
predecessors.  In particular, this description does not charge for
wires or impose any restriction on the support of a linear form.

Suppose first that every $e_j<t_0$.  By
\eqref{eq:low-degree-distances},
\[
 \sum_j d_{e_j}\ge\frac{s}{2n}\sum_j e_j=\frac s2.
\]
Apply Lemmas~\ref{lem:product-rank} and~\ref{lem:degree-estimates},
weakening $q^{d_{e_j}}$ to $q^{d_{e_j}/2}$ for even degrees.  The
combined even-degree correction is at most $\exp(4C_1\sqrt n)$.
Consequently
\begin{equation}
\label{eq:all-low-product}
 R(G)\le2D\exp(4C_1\sqrt n)q^{s/4}
 \le2D\exp(4C_1\sqrt n)n^{-\sqrt n/20},
\end{equation}
where we used $q\le n^{-2/5}$ and $s\ge\sqrt n/2$.

Now suppose that one factor occurrence has degree $e\ge t_0$.
Expand that occurrence, using its middle sum, into at most $S$
products of $e$ homogeneous linear forms.  All other factors remain
unexpanded.  This includes any other occurrences of the same gate:
expanding one occurrence in $Q^h$ leaves its other $h-1$ occurrences
unexpanded.

In each resulting product, the $e$ newly supplied linear factors
contribute $H_1^e\le q^{\lambda e/2}$ to the estimate of
Lemma~\ref{lem:product-rank}, because $d_1=\lambda$ and $1<t_0$.
There are at most $n/t_0=4s$ remaining factor occurrences of degree
at least $t_0$, and their $H$-factors are at most 5 each.  The
remaining low odd degrees contribute at most 1.  The remaining low
even degrees contribute at most $\exp(4C_1\sqrt n)$ after their
favorable powers of $q$ are dropped.  Rank subadditivity over the
expansion therefore yields
\begin{align}
\label{eq:one-high-product}
 R(G)
 &\le2SD\,5^{4s}\exp(4C_1\sqrt n)q^{\lambda t_0/2}\notag\\
 &\le2SD\,5^{4s}\exp(4C_1\sqrt n)n^{-3\sqrt n/160}.
\end{align}
For the last step, $q\le n^{-2/5}$ and
$\lambda\ge3/8$, $t_0\ge\sqrt n/4$ imply
$\lambda t_0/5\ge3\sqrt n/160$.

Since $5^{4s}\le\exp(4\log(5)\sqrt n)$, both
\eqref{eq:all-low-product} and~\eqref{eq:one-high-product} are at most
\[
 2SD\exp(C\sqrt n)n^{-\sqrt n/100}
\]
for an absolute $C$.  Multiplication by a nonzero scalar leaves rank
unchanged, and multiplication by zero gives rank zero.  Summing the
at most $S$ upper-product contributions proves the proposition.
\end{proof}
